% BHU PYQ -- Professional Exam Template
\documentclass[11pt,a4paper]{article}

\usepackage[a4paper,top=2.5cm,bottom=2.5cm,left=2.8cm,right=2.8cm,headheight=14pt]{geometry}
\usepackage{amsmath,amssymb}
\usepackage[T1]{fontenc}
\usepackage[utf8]{inputenc}
\usepackage{lmodern}
\usepackage{microtype}
\usepackage{fancyhdr}
\usepackage{enumitem}
\usepackage{listings}
\usepackage{hyperref}
\usepackage{lastpage}
\usepackage{parskip}

\hypersetup{colorlinks=true,urlcolor=black,linkcolor=black,
  pdftitle={CSB-401: Numerical Computing -- BHU PYQ}}

\pagestyle{fancy}
\fancyhf{}
\renewcommand{\headrulewidth}{0.4pt}
\renewcommand{\footrulewidth}{0.4pt}
\fancyhead[L]{\small   \textit{Banaras Hindu University}}
\fancyhead[R]{\small   \textit{CSB-401: Numerical Computing}}
\fancyfoot[C]{\small Page  \thepage\ of \pageref{LastPage}}

\lstset{
  basicstyle=\small tfamily,
  frame=single,
  breaklines=true,
  columns=flexible,
  keepspaces=true
}


\newlist{parts}{enumerate}{2}
\setlist[parts,1]{label=(\alph*),leftmargin=*,itemsep=4pt,topsep=3pt}
\setlist[parts,2]{label=(\roman*),leftmargin=*,itemsep=2pt,topsep=2pt}

\newcommand{\pts}[1]{\hfill   \textit{[#1]}}


\begin{document}

\begin{center}
  {\large\bfseries BANARAS HINDU UNIVERSITY}\\[2pt]
  {\normalsize B.Sc. (Hons.) Semester IV Examination, 2014}\\[6pt]
  \rule{0.8\linewidth}{0.5pt}\\[6pt]
  {\large\bfseries Paper: CSB-401 --- Numerical Computing}\\[4pt]
  {\normalsize Time: Three hours \hfill Maximum Marks: 70}
\end{center}

\rule{\linewidth}{0.4pt}

\smallskip



\noindent   \textit{Instructions: Answer five questions including Question No. 1 which is compulsory. Use of calculator is allowed. All undefined symbols have their usual meanings.}

\bigskip

\medskip


\noindent   \textbf{Question 1.}

\begin{parts}
  \item Find the maximum absolute error and maximum relative error in $3x^2 y / z$ when $\delta x = \delta y = \delta z = 0.001$ at $x = y = z = 1$. \pts{3}
  \item Define Normalized floating point number. Represent $54.75  \times 10^7$ in normalized floating point mode. \pts{3}
  \item Prove or disprove the relations: \pts{4}
  
\begin{parts}
    \item $\Delta^n \equiv \delta^n E^{n/2}$
    \item $\Delta + 
abla \equiv \frac{\Delta}{
abla} - \frac{
abla}{\Delta}$
  \end{parts}
  \item Find the missing value of the following data: \pts{3}
  
\begin{center}
    
\begin{tabular}{|c|c|c|c|c|c|}
      \hline
      $x$ & 1 & 2 & 3 & 4 & 5 \\
      \hline
      $y$ & 7 & -- & 13 & 21 & 31 \\
      \hline
    \end{tabular}
  \end{center}
  \item Using Jacobi method, find all the eigenvalues and corresponding eigenvectors of the matrix: \pts{3}
  \[ A = 
\begin{bmatrix} 3 & \sqrt{3} \\ \sqrt{3} & 3 \end{bmatrix} \]
  \item Using forward differences, obtain the formula for the second derivative of $y(x)$ at $x = x_0$. \pts{3}
  \item Use Runge-Kutta second order method to approximate $y$ when $x = 0.1$, given that $\frac{dy}{dx} = x + y$, $y(0) = 1$. \pts{3}
\end{parts}

\medskip


\noindent   \textbf{Question 2.}

\begin{parts}
  \item Solve the equation $x  an x + 1 = 0$ by Regula-Falsi method starting with $x_0 = 2.5$ and $x_1 = 3$ correct to three decimal places. \pts{6}
  \item Perform four iterations of Gauss-Seidel method to solve the equations: \pts{6}
  \[ x + 2y + 10z = 61 \]
  \[ 10x + y + 2z = 44 \]
  \[ 2x + 10y + z = 51 \]
  by taking the initial approximation as $[0, 0, 0]$.
\end{parts}

\medskip


\noindent   \textbf{Question 3.}

\begin{parts}
  \item Obtain the order of convergence of the secant method in finding the root of the equation $f(x) = 0$. \pts{6}
  \item From the following table of values of $x$ and $y$, obtain $\frac{d^2y}{dx^2}$ at $x = 2.2$: \pts{6}
  
\begin{center}
    
\begin{tabular}{|c|c|c|c|c|c|c|c|}
      \hline
      $x$ & 1.0 & 1.2 & 1.4 & 1.6 & 1.8 & 2.0 & 2.2 \\
      \hline
      $y$ & 2.7183 & 3.3201 & 4.0552 & 4.9530 & 6.0496 & 7.3891 & 9.0250 \\
      \hline
    \end{tabular}
  \end{center}
\end{parts}

\medskip


\noindent   \textbf{Question 4.}

\begin{parts}
  \item Using the iterative method, find the inverse of the matrix
  \[ A = 
\begin{bmatrix} 1 & 10 & 1 \\ 2 & 0 & 1 \\ 3 & 3 & 2 \end{bmatrix} \]
  by assuming the initial inverse of the matrix to be
  \[ 
\begin{bmatrix} 0.4 & 2.4 & -1.4 \\ 0.14 & 0.14 & -0.14 \\ -0.85 & -3.8 & 2.8 \end{bmatrix} \] \pts{6}
  \item By means of Newton's divided difference formula, find a polynomial and hence find the value of $f(8)$ from the following data: \pts{6}
  
\begin{center}
    
\begin{tabular}{|c|c|c|c|c|c|c|}
      \hline
      $x$ & 4 & 5 & 7 & 10 & 11 & 13 \\
      \hline
      $y$ & 48 & 100 & 294 & 900 & 1210 & 2028 \\
      \hline
    \end{tabular}
  \end{center}
\end{parts}

\medskip


\noindent   \textbf{Question 5.}

\begin{parts}
  \item Prove the identity: \pts{6}
  \[ y_x - \Delta^2 y_x + \Delta^3 y_x - \Delta^5 y_x + \Delta^6 y_x - \Delta^8 y_x + \dots = y_x - \Delta^2 y_{x-1} + \Delta^4 y_{x-2} - \Delta^6 y_{x-3} + \Delta^8 y_{x-4} - \dots \]
  \item Solve the following differential equation using Runge-Kutta method of order four: \pts{6}
  \[ \frac{dy}{dx} = xy + y^2, \quad y(0) = 1, \]
  at $x = 0.1$ and $0.2$.
\end{parts}

\medskip


\noindent   \textbf{Question 6.}

\begin{parts}
  \item Derive Newton-Cotes quadrature formula. Hence obtain Simpson's one-third rule of integration. \pts{6}
  \item Apply Gauss-Jordan method to solve the system of linear equations: \pts{6}
  \[ 2x + 3y - z = 5 \]
  \[ 4x + 4y - 3z = 3 \]
  \[ 2x - 3y + 2z = 2 \]
\end{parts}

\medskip


\noindent   \textbf{Question 7.}

\begin{parts}
  \item Using Milne's method, find $y(0.8)$ if $y(x)$ is the solution of $\frac{dy}{dx} = 1 - y^2$, given that $y(0) = 0, y(0.2) = 0.2027, y(0.4) = 0.4228, y(0.6) = 0.6844$. \pts{6}
  \item Derive two-point Gauss-Legendre formula for integration. \pts{6}
\end{parts}

\vfill
\rule{\linewidth}{0.4pt}

\begin{center}\small
\href{https://bhu-exam.vercel.app/}{Click here to visit our website}\\[4pt]\href{https://me-aryan.vercel.app/}{Click here to visit our contributor}\\[4pt]\href{https://quashan-me.vercel.app/}{Click here to visit our contributor}
\end{center}
\end{document}